% frontmatter/abstrak.tex -- ABSTRAK

\chapter*{ABSTRAK}
\addcontentsline{toc}{chapter}{ABSTRAK}
Ulasan produk berbahasa Indonesia pada platform perdagangan
elektronik terus bertambah setiap hari, sehingga penilaian opini
konsumen secara manual menjadi tidak efisien. Penelitian ini
bertujuan merancang, membangun, dan mengevaluasi model klasifikasi
sentimen berbasis Long Short-Term Memory (LSTM) untuk
mengklasifikasikan ulasan produk ke dalam tiga kelas, yaitu positif,
negatif, dan netral, serta membandingkan kinerjanya dengan metode
pembelajaran mesin konvensional berbasis Support Vector Machine
(SVM) sebagai \textit{baseline}. Metode yang digunakan meliputi
praproses data yang disesuaikan dengan karakteristik morfologis
bahasa Indonesia, seperti normalisasi kata tidak baku dan penghapusan
stopword, dilanjutkan dengan ekstraksi fitur menggunakan
\textit{word embedding} dan pelatihan model LSTM dengan satu lapisan
tersembunyi berisi 128 unit. Kinerja model dievaluasi menggunakan
metrik akurasi, presisi, recall, dan F1-score pada data uji yang
terpisah dari data latih. Hasil pengujian menunjukkan bahwa model
LSTM yang diusulkan memperoleh F1-score sebesar 0,85, lebih tinggi
dibandingkan model \textit{baseline} SVM yang hanya memperoleh
F1-score sebesar 0,75. Peningkatan kinerja tersebut terutama terlihat
pada kemampuan model LSTM dalam menangani kalimat yang mengandung
pola negasi, meskipun model masih menunjukkan keterbatasan dalam
menangani ulasan yang mengandung sarkasme atau opini implisit. Hasil
penelitian ini diharapkan dapat dimanfaatkan sebagai dasar
pengembangan sistem pemantauan opini konsumen secara otomatis.

\vskip 2\baselineskip
\noindent\textbf{Kata kunci : analisis sentimen, LSTM, ulasan produk,
Support Vector Machine}
